\section{Zero momentum and instantaneous quartic transport}
The distinction between momentum and quartic mass transport is exact.
Choose a real sharp wave $Q$ and a real smooth test $h$ with
\[
 \langle h,Q\rangle=0,\qquad
 \langle h,\partial_kQ\rangle=0\ (1\le k\le N),\qquad
 \langle h,\partial_j\Delta Q\rangle\ne0
\]
for some $j$. Such a test exists. Otherwise
$\partial_j\Delta Q\in\operatorname{span}\{Q,\partial_1Q,\ldots,
\partial_NQ\}$ as distributions. Pairing with $Q$ removes its $Q$
component, and Fourier transformation would force $\widehat Q$ onto
the zero set of a nonzero polynomial, contradicting $Q\ne0$ in $L^2$.
For fixed sufficiently large $a>0$, set
\[
 u_\varepsilon=(1+a\varepsilon^2)Q+i\varepsilon h.
\]
Then $P(u_\varepsilon)=0$ exactly, and
\begin{align*}
 d_\varepsilon
 &=\frac4N\left(2a+\frac{M(h)}{M_c}\right)\varepsilon^2
       +O(\varepsilon^4),\\
 E(u_\varepsilon)
 &=\varepsilon^2\left[
 \frac{A(h)-\int|Q|^{8/N}h^2}{2}-a\omega M_c\right]
       +o(\varepsilon^2)<0,
\end{align*}
where $\Delta^2Q+\omega Q-|Q|^{8/N}Q=0$.
Nevertheless its quartic mass current is
\[
 4\int\xi_j|\xi|^2|\widehat u_\varepsilon(\xi)|^2\,d\xi
 =8\varepsilon(1+a\varepsilon^2)
             \langle h,\partial_j\Delta Q\rangle.
\]
Thus this current has order $\sqrt{d_\varepsilon}$ even at exactly zero
momentum and negative energy. With a finite first moment it is the
instantaneous derivative of that moment; without one it remains the
finite Fourier current underlying localized mass transport.


The improved center in Section~\ref{sec:r13-center} is a modulation
coordinate chosen by dual orthogonality, not the ordinary first-moment
center. Its quadratic drift is therefore consistent with the example.
The example rules out the proposed instantaneous current improvement;
it does not rule out further cancellation along the nonlinear evolution.
